%========================================================
\chapter{Conclusion}
%========================================================

This study presents an efficient implementation of the scalable and secure fed
erated graph learning framework for the intelligent TBML detection within distribu ted institutional environments. The designed framework included hetero geneous graph learning, federated Hybrid HGNN-LSTM analytical inference, Apache Kafka-based transaction processing, PostgreSQL-based compli ance orchestration, and regulator-oriented monitoring synchronization within the consoli dated infrastructure. Stable performance in terms of the Fraud-class Precision near 0.98 and reliable federated learning behavior was demonstrated through the experimenta l evaluation of the proposed framework. In addition, concurrent scalability experiments
proved stability in dashboard latency near 156-162 ms and fraud-filtering latenc y near 148-159 ms under gradually increasing monitoring load. All these outcomes pro ve the ability of the proposed framework to perform TBML detection within distri buted institutional environments reliably, scalably, and securely.

%========================================================
\section{Limitations}
%========================================================
In spite of stable analytical performance and effective monitoring capability, there are
some practical limitations associated with the current implementation of the proposed
framework. In particular, experimental results were obtained based on the synthetic
enrichment of institutional transactional datasets due to the lack of institution ally verified TBML datasets annotated with laundering activities. As a consequence, some laundering behaviors and trade manipulation practices might not be sufficiently repre sented within the present experimental environment.

 The federated workflow also results in increased communication costs while
synchronizing parameters in a distributed manner. The synchronization costs may be pro
portionally higher with an increasing number of participating banking environments. In
the same way, the graph-processing workflow involving large banking transaction envi
ronments may cause extra computational costs in the context of high-throughput mul
ti-hop graph traversal and distributed inference operations.
Fraud detection, monitoring, and compliance notification are the primary applica
tions covered by our current system. Although the framework provides effective support
for alert propagation and coordination within regulatory institutions and among institu
tional participants, no automated action mechanisms (such as transaction blocking,
account freezing, document locking, and dynamic compliance triggers) are considered in
the framework, requiring closer coupling with banking infrastructure.

%========================================================
\section{Future Scope}
%========================================================
There are still several possibilities for further scaling up the suggested framework to more production-ready and scalable financial crime-monitoring ecosystems. A future research area would be to test the framework with more institutionally-backed TBML datasets, with real-world laundering annotations, cross-border laundering behaviors and increasingly changing topologies in the financial crime domain. These datasets will enhance the robustness of the analytical and the operational generalisation in large scale banking environments.

In further extensions, the semantic representation extraction system from unstructured trade documents such as invoices, customs declarations, shipping documents, and letters of credit may be implemented using a transformer architecture. These context-specific embeddings can then be combined into the graph-learning process to enhance semantic anomaly detection and trade-document authentication steps. The federated learning ecosystem could also expand to geographically distributed banking ecosystems with more widespread and larger institutional participation.

Future implementations may also utilize enhanced secure aggregation methods like Fully Homomorphic Encryption (FHE), Secure Multi-Party Computation (SMPC), and differential privacy methods to further bolster the security of distributed parameter synchronization operations against gradient-leakage and model-inversion attacks. There is also potential for further research on Continuous-Time Dynamic Graph (CTDG) architectures, in addition to the use of explainable graph neural networks (XGNNs) to further enhance the temporal topology tracking and graph-level interpretability for institutional compliance analysts.

This proposed solution can also be developed into an end-to-end banking compliance solution. In addition to fraud detection and propagation of notifications, future deployments can enable automated compliance-response orchestration, such as freezing transactions, executing risk triggers as per the regulators' requirements, verifying documents from institutions, implementing regulator-approved intervention pipelines, and restricting accounts based on policies embedded in core banking systems.
%========================================================
\section{Learning Outcomes of the Project}
%========================================================

% \begin{itemize}

% \item Acquired practical experience in designing and deploying heterogeneous graph-learning architectures for large-scale financial fraud-detection workflows.

% \item Developed understanding of federated learning environments, distributed parameter synchronization, and privacy-preserving collaborative analytical inference across institutional systems.

% \item Gained hands-on exposure to real-time streaming architectures using Apache Kafka, asynchronous backend synchronization, and distributed event-driven monitoring workflows.

% \item Learned implementation and optimization of graph neural networks, temporal sequence modeling, and multi-modal feature fusion for fraud-classification tasks.

% \item Developed expertise in PostgreSQL database optimization, concurrent API scalability analysis, connection pooling, and distributed compliance-management synchronization workflows.

% \item Acquired experience in full-stack system integration involving Next.js dashboards, Prisma ORM synchronization, authentication workflows, and regulator-facing monitoring environments.

% \item Gained understanding of operational AML workflows including STR lifecycle management, compliance orchestration, fraud-notification propagation, and institutional investigation coordination.

% \item Learned experimental benchmarking, model-performance evaluation, Precision--Recall analysis, confusion-matrix interpretation, and distributed scalability validation for real-time analytical systems.

% \end{itemize}

The suggested TBML detection system introduced practical experience with het
erogenous graph-learning models, federated learning platforms, and real-time financial
fraud-detection systems. The project gave practical experience working with Graph Neural Net
works, time sequence analysis, synchronized model parameters training, and data infer
ence using privacy-preserving analytics. Moreover, practical knowledge was obtained re
lated to Apache Kafka real-time stream processing, PostgreSQL database manage
ment, scalability evaluation of an API, and end-to-end application development using
Next.js, Prisma ORMs, and FastAPI. Also, practical skills were developed concerning
AML compliance processes, STR lifecycle, fraud detection monitoring, precision
recall analysis, and distributed scalability testing in real-time systems.

